\section{General Overview}
\label{sec:general_overview}

% Describe the functionalities your system will provide from the end user's perspective.
% Present an overview of your system, outlining its key software and hardware components 
% without going into implementation details.
% Imagine a real full-scale deployment of your solution.

% TODO: Write general overview

\subsection{System Functionalities}
% Describe functionalities from the end user's perspective
Users can monitor their sleep quality and access sleep-enhancing
features via bone conduction sound waves. The system displays real-time SpO2​
levels and heart rate data on a separate device. Since the user is a patient within a facility,
their doctor also has access to the system. Through a dedicated interface, the doctor can
monitor all patient stats live and remotely adjust the sound wave frequencies to optimize
sleep improvement.
\subsection{Mobile Components}
% The mobile component(s) and their role in the solution
The system features a sleep-monitoring component in the form of an elastic headband worn by the patient.
This band monitors heart rate and SpO2​ while simultaneously inducing sleep-enhancing sound waves.
When the patient goes home, they use a separate device equipped with an LCD screen to display their
vitals. This second device acts as a gateway; the user simply uses their phone to configure the initial
Wi-Fi connection.
\subsection{Server and Cloud Services}
% The server or cloud services functionalities needed
The infrastructure relies on an MQTT server to ingest real-time statistics from each patient.
Data is stored in a time-series database for long-term tracking, while a dedicated
backend and frontend power the application used by doctors to monitor and manage their patients.
\subsection{IoT and MCU Devices}
% Types of IoT or MCU devices, their sensors and actuators
The sleep-tracking device is powered by an ESP32 integrated with two primary sensors:
an MPU-6050 (featuring a 3-axis gyroscope and a 3-axis accelerometer) to monitor patient movement and
sleep states, and a MAX30102 to track SpO2​ and heart rate. A bone conduction speaker serves as the
primary actuator for sleep induction. The secondary home device also utilizes an ESP32 and features an
OLED display as an actuator to visualize the patient's vital signs.
\subsection{Component Interactions}
% What components will need to interact/communicate
The sleep-tracking headband connects to the gateway device via \textbf{Wi-Fi Direct}
to transmit real-time sensor data. For initial setup, the gateway device pairs with the patient's
smartphone via \textbf{Bluetooth} to configure local network credentials. Once connected to the internet,
the gateway uses the \textbf{MQTT} protocol to relay the collected sleep statistics to the central server.
This creates a bridge where the tracking device sends data to the gateway, which then pushes it to
the cloud, allowing the doctor's frontend to access the information and send frequency adjustments
back through the same pipeline.

\subsection{Demonstration Prototype}
% How many of each component you propose to implement for the demonstration prototype
While we plan to implement the full application, there are a few hardware limitations to note.
Due to time constraints, the devices will not yet include integrated batteries. Additionally,
custom 3D-printed plastic cases still need to be developed to house the components and make the
devices wearable and user-friendly for clients.
